\documentclass[12pt]{article}
\usepackage[utf8]{inputenc}
\usepackage[T1]{fontenc}
\usepackage{hyperref}
\title{Tuotomb Language Resources for K-12 Curriculum Development}
\author{Compiled by Claude Sonnet 5 (claude-sonnet-5) for Andr\'as Kornai}
\date{}
\begin{document}
\maketitle
\section{Introduction}
Tuotomb (ISO 639-3 ttf; also called Bonek, Otomb, Atomb) is a Southern Bantoid language of Cameroon, centred on the village of Bonek near Ndikinemeki, with about 1,000 speakers (2007). It is classified as endangered. Very little published material exists.
\section{Primary Educational Resources}
No primers, textbooks or school materials were found in search results. A 2025 revitalisation project listed on the Endangered Languages Project reportedly aims to produce a community orthography, an online dictionary, school materials and recordings; the outputs were not verified.
\begin{itemize}
\item Endangered Languages Project entry (project description): \url{https://endangeredlanguages.com/node/114877}
\end{itemize}
\section{Linguistic and Reference Materials}
\begin{itemize}
\item Wikipedia, Tuotomb language: \url{https://en.wikipedia.org/wiki/Tuotomb_language}
\item Numeral system survey (MPI-EVA), citing fieldwork by G. Boyd, 2009: \url{https://lingweb.eva.mpg.de/channumerals/Tuotomb.htm}
\end{itemize}
\section{Religious and Community Resources}
Joshua Project reports no Bible translation or scripture resources.
\begin{itemize}
\item Joshua Project language page: \url{https://joshuaproject.net/languages/ttf}
\item Joshua Project people group page: \url{https://joshuaproject.net/index.php/people_groups/15646}
\end{itemize}
\section{Digital and Audio Resources}
No audio, video or corpus resources were found in search results.
\section{Curriculum Development Context}
The language is not known to be taught in schools and has no established orthography in the sources found. Any curriculum would have to start from community-led orthography and documentation work, possibly building on the project above.
\end{document}